\needspace{12\baselineskip}
\begin{center}
\pdfbookmark[1]{8. Zu Art. 114 und 177.}{zusatz-008}
[8. Zu Art. 114 und 177.]
\end{center}
\begin{center}
\{Druckfehler in Dr. \textsc{Gauss}' Theoria motus corporum coelestium etc. Hamburgi 1809, von Senator Bar. \textsc{Oriani}.\}
\end{center}
\begin{center}
\begin{tabular}{c}\hline
\small Monatliche Correspondenz Band XXI, S. 282--283, März 1810.\\\hline
\end{tabular}
\end{center}
\{.... Damit Gleichung 7 [Art. 114] eine Identität wird, müssen die Koeffizienten von \(\delta\delta'\delta''\), \(D\delta'\delta''\), \(DD'\delta''\) usw. \(=0\) sein; nach den gehörigen Reduktionen ist aber der Koeffizient von \(\delta\delta'\delta''\)
\SourceLeaf{308}
\[
(0.1.2)=\left\{\begin{aligned}
&\operatorname{tang}\beta'\operatorname{tang}\beta\,\sin(L''-\alpha'')\sin(L'-L)\\
&+\operatorname{tang}\beta'\operatorname{tang}B\,\sin(L''-\alpha'')\sin(L'-\alpha)\\
&+\operatorname{tang}\beta\operatorname{tang}\beta''\sin(L'-\alpha')\sin(L-L'')\\
&+\operatorname{tang}B\operatorname{tang}\beta''\sin(L'-\alpha')\sin(\alpha-L'')\\
&+\operatorname{tang}\beta''\operatorname{tang}\beta'\sin(L-\alpha)\sin(L''-L')\\
&+\operatorname{tang}\beta''\operatorname{tang}B'\sin(L-\alpha)\sin(L''-\alpha')\\
&+\operatorname{tang}B'\operatorname{tang}\beta\,\sin(L''-\alpha'')\sin(\alpha'-L)\\
&+\operatorname{tang}B'\operatorname{tang}B\,\sin(L''-\alpha'')\sin(\alpha'-\alpha)\\
&+\operatorname{tang}\beta\operatorname{tang}B''\sin(L'-\alpha')\sin(L-\alpha'')\\
&+\operatorname{tang}B\operatorname{tang}B''\sin(L'-\alpha')\sin(\alpha-\alpha'')\\
&+\operatorname{tang}B''\operatorname{tang}\beta'\sin(L-\alpha)\sin(\alpha''-L')\\
&+\operatorname{tang}B''\operatorname{tang}B'\sin(L-\alpha)\sin(\alpha''-\alpha').
\end{aligned}\right.
\]
Gilt daher lediglich
\[
\left.\begin{aligned}
&\operatorname{tang}\beta'\operatorname{tang}\beta\,\sin(L''-\alpha'')\sin(L'-L)\\
&+\operatorname{tang}\beta\operatorname{tang}\beta''\sin(L'-\alpha')\sin(L-L'')\\
&+\operatorname{tang}\beta''\operatorname{tang}\beta'\sin(L-\alpha)\sin(L''-L')
\end{aligned}\right\}=0,
\]
so wird derselbe Koeffizient nicht \(=0\), sondern es ist außerdem erforderlich, dass
\[
B=0,\qquad B'=0,\qquad B''=0.\}
\]
\begin{center}\rule{.25\textwidth}{.35pt}\end{center}

\{Der elegante Satz, der [Art. 117] dem berühmten \textsc{Laplace} zugeschrieben wird, ist tatsächlich zuerst von \textsc{Leonardo Eulero} gefunden worden. Denn in Comment. Acad. Petropol. Tom. XVI\HistoricalNoteEight{Handschriftliche Bemerkung von \textsc{Gauss}:] \textls[150]{Dieser} Satz findet sich an der angegebenen Stelle nicht, wohl aber auf S. 101 Folgendes: \(\int_0^1\BareRoot{\left(\log\frac1x\right)}\mathrm{d}x=\tfrac12\BareRoot{\pi}.\) Setzt man hier \(x=e^{-tt}\), so wird \(\int_0^\infty2tte^{-tt}\mathrm{d}t=\tfrac12\BareRoot{\pi}.\) Es ist aber \(2tte^{-tt}\mathrm{d}t=-\mathrm{d}(te^{-tt})+e^{-tt}\mathrm{d}t\) und \(te^{-tt}=0\) sowohl für \(t=0\) als auch für \(t=\infty\), also \(\int_0^\infty e^{-tt}\mathrm{d}t=\tfrac12\BareRoot{\pi}.\)} zeigt \textsc{Eulerus}, dass das Integral \(-\int\frac{\mathrm{d}x}{\BareRoot{\left(\log\frac1x\right)}}\), von \(x=1\) bis \(x=0\) genommen, \(=\BareRoot{\pi}\) ist, wobei \(\pi\) der halbe Umfang eines mit dem Radius \(=1\) gezeichneten Kreises ist. Setzt man nun \(x=e^{-tt}\), so erhält man
\[
\frac{-\mathrm{d}x}{\BareRoot{\left(\log\frac1x\right)}}=2e^{-tt}\mathrm{d}t.
\]

Daher ist das Integral \(\int e^{-tt}\mathrm{d}t\) von \(t=0\) bis \(t=\infty\) \(=\tfrac12\BareRoot{\pi}\), und folglich wird dasselbe Integral von \(t=-\infty\) bis \(t=+\infty\) \(=\BareRoot{\pi}.\)\}
\begin{center}\rule{.25\textwidth}{.35pt}\end{center}

\SourceLeaf{309}
\begin{center}
\begin{tabular}{c}\hline
\small Monatliche Correspondenz. Band XXI. S. 280. März 1810.\\\hline
\end{tabular}
\end{center}
Für die Notierung der Druckfehler in meiner Theoria bin ich Herrn \textsc{Oriani} sehr verbunden. Er hat ganz recht, dass ich auf S. 129 [Art. 114] hinzuzufügen vergessen habe, dass \(B,B',B''=0\) vorausgesetzt werden müssen, wenn die Bedingungsgleichung, bei der Gleichung 7 unbrauchbar ist, die dort angegebene Gestalt haben soll. Es ist übrigens klar, dass auch dann, wenn nicht \(B,B',B''=0\) gilt, Gleichung 7 unbrauchbar sein kann, nämlich wenn der 12-gliedrige Ausdruck, den \textsc{Oriani} entwickelt hat, zufällig \(=0\) oder sehr klein wird.

Dass \textsc{Euler} schon den Satz gefunden hat, aus dem der schöne, von mir \textsc{Laplace} zugeschriebene Lehrsatz sehr leicht abgeleitet werden kann, fiel mir selbst schon früher ein, allerdings erst, als die Stelle auf S. 212 [Art. 177] bereits gedruckt war; ich wollte es jedoch nicht unter die Errata setzen, weil \textsc{Laplace} wenigstens den obigen Satz zuerst in der dort verwendeten Form aufgestellt hat. ....
\begin{center}\rule{.18\textwidth}{.35pt}\end{center}
